\documentclass[11pt]{article}
\usepackage{mathblatt}
% Prüfungsblatt nach mathe-nachhilfe pruefung.md § 5 (Runde Dreiecke, 11.10.2026), Hand-Satz; Präambel des Setzers.
% Präambel des Setzers (werkzeuge/setzer.py), wortgleich aus dem Hand-Satz T6B (bau/proben/2026-10-09/kathete-berechnen), dort aus K4W.
\makeatletter
\setlength{\mbpfbe}{0mm}% keine Punkte (Bauregel 6.4)
% eingerückt (Bauregel 4.7, 6.6): \gEin Einzug, \gSchrift eine Stufe kleiner
\newlength{\gEin}\setlength{\gEin}{0pt}
\let\gSchrift\relax
\renewcommand{\pfkreuzzeile}[1]{\par\noindent\hangindent3.2em\hangafter1\hspace*{1.6em}$\square$\ #1\par}
\newcommand{\gkreuz}[1]{$\square$\ #1\hspace{2em}}
% Bauregel 6.7: links, was man ansieht, rechts, was man tut; Spaltengrenze überall an derselben Stelle
\newsavebox{\gbox}\newlength{\gLinks}
\newcommand{\gzwei}[2]{\par\smallskip\noindent\sbox{\gbox}{#1}%
  \setlength{\gLinks}{\dimexpr0.5\textwidth-\gEin-\mbpfnr\relax}%
  \ifdim\wd\gbox>\dimexpr\gLinks-4mm\relax
    \usebox{\gbox}\par\smallskip\noindent\begin{minipage}{\linewidth}#2\end{minipage}%
  \else
    \begin{minipage}[t]{\gLinks}\vspace{0pt}\usebox{\gbox}\end{minipage}%
    \begin{minipage}[t]{\dimexpr\linewidth-\gLinks\relax}\vspace{0pt}\raggedright #2\end{minipage}%
  \fi\par}
\RenewDocumentEnvironment{pfaufg}{m m m}{%
  \begin{lrbox}{\mbaufgabebox}\begin{minipage}[t]{\linewidth}%
  \gSchrift\noindent\hspace*{\gEin}\mbpfrandmarke{#1}%
  \begin{minipage}[t]{\mbpfnr}\strut\textbf{#2}\end{minipage}%
  \begin{minipage}[t]{\dimexpr\linewidth-\gEin-\mbpfnr\relax}\raggedright\strut\ignorespaces}%
 {\end{minipage}%
  \end{minipage}\end{lrbox}%
  \setlength{\mbaufgabehoehe}{\dimexpr\ht\mbaufgabebox+\dp\mbaufgabebox\relax}%
  \par\addvspace{7pt}\Needspace*{\mbaufgabehoehe}%
  \noindent\usebox{\mbaufgabebox}\par\addvspace{7pt}}
\makeatother
% Rechenraum als Karo (Bauregel 6.9): n Rechenschritte = n mal 10 mm
\newcommand{\karo}[1]{\par\smallskip\noindent\begin{tikzpicture}
  \clip (0,0) rectangle ({\linewidth-1mm},{#1*10mm});
  \draw[black!16,line width=0.3pt,step=5mm] (0,0) grid ({\linewidth},{#1*10mm});
  \end{tikzpicture}\par}
\newcommand{\kk}[1]{$\square$\,#1\hspace{0.9em}}
\newcommand{\linie}[1]{\,{\color{black!45}\rule[-1pt]{#1}{0.4pt}}\,}
% Zeile am Ende jedes Durchgangs: Originale klein grau, darüber; dann Vorgänger/Nachfolger (Bauregel 3.4, 6.5)
\newcommand{\blattende}[2]{\par\vfill\noindent\begin{minipage}{\linewidth}{\scriptsize\color{mbgrau}Originale: #1\par}\smallskip
{\footnotesize #2\par}\end{minipage}}
\newcommand{\pkt}{\textperiodcentered{}}

\newcommand{\kennung}{NUB}
\newcommand{\figbox}[2]{\begin{minipage}[b]{0.24\linewidth}\centering{\footnotesize #1}\par\smallskip #2\par\smallskip\linie{10mm}\ Achsen\end{minipage}}

\begin{document}
\pfheftstil
\pfheftkopf[\kennung]{Symmetrie}{P10}

\Needspace*{0.25\textheight}\pfstufekopf{Symmetrieachsen zählen}{geprüft 2025 \pkt{} 2022}

% Bank symmetrie-abbildungen-e2-k2-s2-v3, s3-v1, s3-v3, s4-v3 (Päckchen: eine, zwei, drei, keine Achse)
\begin{pfaufg}{}{1.}{}
Wie viele Symmetrieachsen hat das Dreieck oder das Viereck?\par\medskip\noindent
\figbox{a) gleichschenklig}{\begin{tikzpicture}[line width=0.7pt,scale=0.42]\draw (0,0) -- (3,0) -- (1.5,3.6) -- cycle;\end{tikzpicture}}\hfill
\figbox{b) Raute}{\begin{tikzpicture}[line width=0.7pt,scale=0.42]\draw (0,1.6) -- (2.2,0) -- (4.4,1.6) -- (2.2,3.2) -- cycle;\end{tikzpicture}}\hfill
\figbox{c) gleichseitig}{\begin{tikzpicture}[line width=0.7pt,scale=0.42]\draw (0,0) -- (3.6,0) -- (1.8,3.118) -- cycle;\end{tikzpicture}}\hfill
\figbox{d) alle Seiten verschieden}{\begin{tikzpicture}[line width=0.7pt,scale=0.42]\draw (0,0) -- (4.4,0) -- (1.2,2.8) -- cycle;\end{tikzpicture}}
\fusshilfe{\mbox{1:~a) 1; b) 2; c) 3; d) 0}}
\end{pfaufg}

% Bank symmetrie-abbildungen-e2-k2-s4-v1 (Fragerichtung umgekehrt, ohne Bild)
\begin{pfaufg}{}{2.}{}
Kreuze an, welches Viereck keine Symmetrieachse hat.\par\smallskip
\pfkreuz{A}{Rechteck}\pfkreuz{B}{Parallelogramm}\pfkreuz{C}{Raute}
\fusshilfe{\mbox{2:~B}}
\end{pfaufg}

% 2022-OS-B1i, herausgelöst (eigener Wortlaut; vier statt sechs Optionen, gemeinsam.md „Satz“)
\begin{pfaufg}{nach~P10\\’22~\pkt{}~1i}{3.}{}
Diese Figur ist ein gleichschenkliges Trapez. Kreuze an, wie viele Symmetrieachsen die Figur hat.
\gzwei{\begin{tikzpicture}[line width=0.7pt,scale=0.9]\filldraw[fill=black!15] (0,0) -- (3,0) -- (2.5,1.6) -- (0.5,1.6) -- cycle;\end{tikzpicture}}{%
\pfkreuz{A}{0}\pfkreuz{B}{1}\pfkreuz{C}{2}\pfkreuz{D}{4}}
\fusshilfe{\mbox{3:~B}}
\end{pfaufg}

% 2025-OS-K2a, herausgelöst (nur die Symmetrieachsen; BF und AB im Blatt „Länge mit Pythagoras“, Maße weggelassen)
\begin{pfaufg}{nach~P10\\’25~\pkt{}~2a}{4.}{}
Die Figur zeigt das Drachenviereck $ABCD$. Seine Diagonalen $\overline{AC}$ und $\overline{BD}$ schneiden sich im Punkt $F$ im rechten Winkel. Kreuze an, wie viele Symmetrieachsen das Drachenviereck hat.
\gzwei{\begin{tikzpicture}[line width=0.7pt,scale=0.042]
\coordinate (A) at (-25,0);\coordinate (C) at (25,0);\coordinate (D) at (0,25);\coordinate (B) at (0,-55);\coordinate (F) at (0,0);
\draw (A) -- (B) -- (C) -- (D) -- cycle;\draw[dashed] (A) -- (C);\draw[dashed] (B) -- (D);\mbrwmarke{(F)}{(C)}{(D)}
\node[font=\small,anchor=east] at (A) {$A$};\node[font=\small,anchor=north] at (B) {$B$};\node[font=\small,anchor=west] at (C) {$C$};\node[font=\small,anchor=south] at (D) {$D$};\node[font=\small,anchor=north east] at (F) {$F$};
\end{tikzpicture}}{%
\pfkreuz{A}{0}\pfkreuz{B}{1}\pfkreuz{C}{2}\pfkreuz{D}{4}}
\fusshilfe{\mbox{4:~B}}
\end{pfaufg}

\pfblattende{11}{26 FOR & 25 & 24 & 23 & 22 & 21 & 20 & 19 & 18 & 17 & 16}%
{ & 2a & & & 1i & 1j & & & & & }%
{Weiter: Flächeninhalt und Umfang}
\end{document}
